% !TeX root = ../main.tex
行文至此，论文暂告一段落。回望求学与研究的这段历程，始终觉得刘禹锡《浪淘沙》中的一句最能映照此时心境：“千淘万漉虽辛苦，吹尽狂沙始到金。”从最初的懵懂摸索，到一次次推翻重来、反复修改，再到今日终于将思考与实践整理成文，这一路并不轻松，也因此更知所得之可贵。

首先，我要衷心感谢我的导师。感谢导师在论文选题、研究思路、技术路线与写作修改过程中给予我的耐心指导与严格要求。每当我在研究中感到迟疑与困惑时，导师总能以严谨的学术态度和宽厚的包容之心，为我指出问题、启发思路。导师不仅教会了我如何开展科研，也让我更加明白踏实治学、认真做事的意义。这份教诲将使我长期受益。

感谢实验室的老师和同学们。与你们共同学习和讨论的日子，是研究生阶段最珍贵的记忆之一。无论是方案交流中的启发、实验推进中的协助，还是日常相处中的鼓励与陪伴，都让我在繁重的学习和科研压力之下感受到温暖与力量。许多看似平常的支持，实际上都成为我坚持走下去的重要原因。

感谢我的家人。你们或许未必完全了解我所研究的具体内容，却始终以最深沉、最无条件的方式支持我。每一次情绪低落时的安慰、每一个重要阶段中的守望，都让我在远行与奔忙之中始终有所依归。正是这份理解、信任与包容，使我能够心无旁骛地完成学业。

同时，我也要感谢一路上给予我关心与帮助的朋友们。感谢你们在我焦虑、疲惫和自我怀疑时的倾听与鼓励，也感谢你们在平凡日常中的真诚陪伴。正因为有你们，这段本该艰难而孤独的旅程，才多了许多明亮与坚定。

最后，感谢过去那个在压力与不确定中仍未轻言放弃的自己。论文的完成并不意味着终点，而是新的起点。愿自己在今后的学习、工作与生活中，仍能保持初心，笃行不怠，不负曾经走过的路，也不负所有真诚的期待。
